Distributed Acoustic Sensing (DAS) menukarkan gentian optik kepada beribu-ribu penderia getaran, dan nilainya bagi keselamatan perimeter bergantung pada keupayaan mengecam peristiwa yang direkodkan oleh penderia tersebut. Rangkaian neural dalam yang mendahului penanda aras DAS awam lazimnya memerlukan unit pemprosesan grafik (GPU), sukar ditafsir, dan biasanya diuji menggunakan pembahagian rawak atau kronologi ke atas tetingkap yang bertindih. Tesis ini bertujuan untuk menganalisis keperluan pengiraan bagi pengelasan peristiwa DAS serta faktor yang menjadikan sesuatu saluran paip ringan, membangunkan saluran paip pembelajaran mesin pelbagai domain bagi tugas tersebut, dan menilainya berbanding garis dasar pembelajaran mendalam yang telah diterbitkan. Berpandukan Design Science Research Methodology (DSRM), saluran paip ini direka bentuk menggunakan set data keselamatan perimeter Tomasov yang mempunyai sembilan kelas, dan digunakan tanpa sebarang perubahan pada set data Cao yang mempunyai enam kelas, yang direkodkan menggunakan sistem penderiaan yang tidak berkaitan. Setiap tetingkap diperihalkan oleh $1{,}002$ deskriptor bernama daripada domain masa, frekuensi, masa--frekuensi dan ruang, yang boleh dikurangkan melalui konsensus lapan kaedah pemilihan ciri, dan tujuh pengelas dibandingkan di bawah enam strategi bagi menangani ketidakseimbangan kelas. Kebaharuan kajian ini bersifat metodologi, bukan algoritmik. Keterasingan data latihan dan ujian ditegaskan dalam kod sebelum sebarang model dilatih, dan saiz, kependaman serta kos latihan setiap model diukur pada papan terbenam yang sama. Audit terhadap protokol yang diterbitkan menunjukkan bahawa, di bawah pembahagian rawak, setiap tetingkap ujian Tomasov berkongsi sampel mentah dengan tetingkap latihan, dan bahawa 423 daripada 441 sesi rakaman Cao terdapat dalam kedua-dua folder yang diterbitkan. Pembahagian berblok dengan jalur pengawal dan pembahagian berkumpulan mengikut sesi menghapuskan kedua-dua bentuk kebocoran ini. Pada pembahagian tersebut, gabungan empat domain mengatasi setiap domain tunggal, dan pengelas LightGBM yang dipilih mencapai skor makro-$F_1$ sebanyak $81.98$ pada Tomasov dan $98.90$ pada Cao. Skor ini melebihi skor convolutional neural networks (CNN) yang diterbitkan bersama kedua-dua set data sebanyak $16.43$ dan $8.01$ mata, serta skor lima seni bina graf dan transformer yang diterbitkan sejak tahun 2024. Pemilihan konsensus membuang kira-kira $95\%$ ciri dengan kehilangan $8.39$ dan $5.81$ mata, manakala strategi pengimbangan kelas memberi kesan yang kecil. Dengan hanya mengubah pembahagian data, ketepatan rangkaian yang sama berubah sehingga $26.27$ mata, dan penghapusan kebocoran menurunkan skor $F_1$ kelas \emph{openclose} daripada $0.72$ yang diterbitkan kepada $0.044$. Keupayaan generalisasi di luar rakaman latihan adalah lebih lemah. Ketepatan purata ialah $63.85\%$ pada rakaman yang tidak pernah dilihat bagi empat kelas Tomasov dan menurun daripada $99.06\%$ kepada $86.81\%$ pada tarikh perolehan Cao yang tidak pernah dilihat, manakala bagi empat jenis peristiwa yang dikongsi oleh kedua-dua set data, model yang dilatih pada satu sistem penderiaan mencatat skor di bawah paras peluang pada sistem yang lain. Pada NVIDIA Jetson AGX Orin, model yang dipilih bersaiz $0.83$--$3.42$\,MB, dilatih menggunakan CPU dalam masa $41$--$1{,}357$\,s, dan mengelaskan satu tetingkap dalam masa kurang daripada $0.2$\,ms, manakala garis dasar rangkaian neural bersaiz sehingga $132.6$\,MB dan memerlukan sehingga sepuluh setengah jam latihan GPU. Pengekstrakan ciri, iaitu $122$--$145$\,ms bagi setiap tetingkap, kekal sebagai peringkat yang paling berat dari segi pengiraan. Dapatan kajian menunjukkan bahawa saluran paip yang ringan dan boleh ditafsir mampu bersaing dengan seni bina pembelajaran mendalam terkini, dan bahawa perbandingan sedemikian hanya bermakna apabila penilaian dikumpulkan mengikut rakaman dan disahkan dalam kod.
